\itemhead{3.1 例题 · 例题6（教材p.74）}{XJ-B1-C03-EX-P074-Y0394}
\label{XJ-B1-C03-EX-P074-Y0394}
{\footnotesize 湘教版教材p.74 / PDF第81页 / 3.1 例题 / 例题6}
\field{原题}{为推进生态优先、节约集约、绿色低碳发展，推动形成绿色低碳的生产方式和生活方式，某地采用分段计费的方法计算用户的电费：每月用电量不超过 $100\,\mathrm{kW\cdot h}$，按 $0.57$ 元/($\mathrm{kW\cdot h}$)计费；每月用电量超过 $100\,\mathrm{kW\cdot h}$，其中 $100\,\mathrm{kW\cdot h}$ 仍按原标准收费，超过部分按 $1.5$ 元/($\mathrm{kW\cdot h}$)计费。
\subq{1}设月用电 $x\,\mathrm{kW\cdot h}$，应交电费 $y$ 元，写出 $y$ 关于 $x$ 的函数解析式。
\subq{2}小赵家第一季度缴纳的电费情况如下表：
\begin{center}\begin{tabular}{c|rrrr}\hline
月份&1&2&3&合计\\
计费金额/元&114&75&45.6&234.6\\\hline
\end{tabular}\end{center}
问：小赵家第一季度共用电多少？}
\field{参考答案}{（1）$y=\begin{cases}0.57x,&0\le x\le100,\\1.5x-93,&x>100.\end{cases}$（2）各月用电138、112、80千瓦时，合计330千瓦时。须把超过100的部分单独计费。}
{\small \field{关联知识}{F07 分段函数、绝对值及分段应用（教材帮p.86、91、92、93 / PDF95、100、101、102）；F03 函数求值、解析式与函数建模（教材帮p.80、87、88、89、90、131 / PDF89、96、97、98、99、140）}
\field{方法与易错点}{先列转折点和各段条件，再建立分段式，检查衔接、定义域与值域。 易错：跨段套公式；超额部分按全量收费；绝对值去号漏区间。}
\field{答案依据}{例题原解位置：湘教版教材p.74及题库原图；独立编写可编辑参考解答。答案册此部分主要收练习、习题，不伪造例题答案册位置。}}
